\section{Введение}

Специфика функционирования малых космических аппаратов (МКА) формата CubeSat на низких околоземных орбитах (LEO) накладывает жесткие ограничения на архитектуру бортовой аппаратуры спутниковой навигации. Высокая орбитальная скорость движения платформы сопровождается значительным доплеровским сдвигом несущей частоты принимаемых радиосигналов GPS L1 и его высокой динамикой. Это многократно расширяет область неопределенности в пространстве параметров «задержка -- Доплер», требуя применения алгоритмов параллельного поиска по кодовой фазе (Parallel Code Phase Search, PCPS) для обеспечения приемлемого времени «холодного» старта навигационного приемника.

Реализация алгоритма PCPS на кристалле связана с непрерывной обработкой потока отсчетов высокой плотности и выполнением ресурсоемких процедур согласованной фильтрации на базе быстрого преобразования Фурье. Традиционный маршрут проектирования таких трактов опирается на проприетарные САПР и событийно-управляемые симуляторы логического уровня. Однако проверка алгоритмов захвата на длинных сигнальных выборках физического эфира сталкивается с непреодолимым вычислительным барьером: моделирование даже нескольких миллисекунд реального радиосигнала требует миллиардов элементарных логических операций, что приводит к критическому замедлению процесса проектирования.

Разрешение указанного противоречия лежит в переходе к методологии совместного аппаратно-программного моделирования (HW/SW Co-Design) на основе открытых инструментов. Применение компилятора Verilator позволяет преобразовать синтезируемое описание вычислительного тракта в тактово-точную программную модель, исключающую накладные расходы классических систем верификации и обеспечивающую бесшовную интеграцию с аналитическими алгоритмами постобработки на C++.

Целью настоящей работы является разработка и экспериментальное исследование сквозного аппаратно-программного комплекса первичного поиска навигационных сигналов GPS L1 C/A. Предложенная методология объединяет синтезируемый RTL-тракт цифровой обработки сигналов и высокопроизводительный верификационный стенд на базе Verilator с адаптивным CFAR-детектором по критерию Неймана--Пирсона, формируя прямой маршрут интеграции в открытые процессорные платформы RISC-V и заказные СБИС.
